\note{2}{Wed 18 Oct 2023 12:33}{Presentation Handout for Polya's Urn}

\section{Polya's Urn Presentation}


\begin{definition}[Polya's Urn Model]
	Let $c>0$ be an integer. The polya's urn model is a Markov process $(R_n, G_n)_{n\ge 0}$ on $\{i \in \mathbb{Z}: i > 0\} ^2$ with transition probabilities
	\begin{align*}
		p\left((i+c, j) \middle| (i,j) \right) &= \frac{i}{i+j}\\
		p\left((i, j+c) \middle| (i,j) \right) &= \frac{j}{i+j}
	.\end{align*}
\end{definition} 

We interpret the tuple as the number of red balls and green balls. Set the initial condition $(R_0, G_0) = (r, g)$. 
\begin{claim}
	The fraction of green balls, $X_n = \frac{G_n}{R_n + G_n}$, is a martingale.
\end{claim}
\begin{proof}
	Given $(R_n, G_n) = (i,j)$, then $X_{n+1}$ has distribution
	\begin{align*}
		P\left(X_{n+1} = \frac{j+c}{i + j + c}\right) &= \frac{j}{i + j}\\
		P\left(X_{n+1} = \frac{j}{i + j + c}\right) &= \frac{i}{i + j}
	.\end{align*}
	Since
	\[
		\frac{j+c}{i+j+c} \cdot \frac{j}{i+j}+\frac{j}{i+j+c} \cdot \frac{i}{i+j}=\frac{(j+c+i) j}{(i+j+c)(i+j)}=\frac{j}{i+j}
	,\] 
	we conclude that
	\[
		E(X_{n+1} | R_n, G_n) = X_n
	.\] 
\end{proof}

Now since $X_n$ is a nonnegative martingale, it converges to a limit $X_\infty $ a.s. To find out its limit, we first compute the distributions of urn state after $n$ draws.

\subsection{Distribution of urn state}

Consider the event that $m$ green balls are drawn consecutively, followed by $(n-m)$ red balls. The probability is
	\[
		\left( \frac{g}{g+r} \cdot \frac{g+c}{g + r + c} \cdot \frac{g + (m - 1) c}{g + r + (m-1) c} \right) 
		\left( \frac{r}{g + r + m c} \cdots \frac{r + (n - m - 1) c}{g + r + (n-1)c}  \right) 
	.\] 
But notice: for any other outcome of first $n$ draws with $m$ green balls has the same probability, since the denominator remains the same and the numerator is permuted.

Use the notation
\[
	r^{(s, j)}=r(r+s)(r+2 s) \cdots[r+(j-1) s]
.\] 
for all $r, s>0$ and natural number $j$. We see that

\begin{lemma}
	\[
	P\left( \left( R_n, G_n \right) = \left( g+ m c, r + (n - m) c \right) \right) 
	= \binom{n}{m} \frac{g^{(c, m)} r^{(c, n-m)}}{(g + r)^{(c, n)}}
.\]
\end{lemma}

Now
\begin{align*}
	\binom{n}{m}\frac{g^{(c, m)} r^{(c, n-m)}}{(g + r)^{(c, n)}}
	&=  \binom{n}{m}\frac{
	\frac{\Gamma\left(g / c+m\right) \Gamma\left(r / c         +n-m\right)}
		{
		\Gamma\left(g / c \right) \Gamma\left(r / c \right)}}
	{
	\frac{\Gamma\left((g + r) / c+n\right)}{\Gamma\left((g + r) / c \right)}} \\
	%&= \frac{B\left(\frac{g}{c}+m, \frac{r}{c}+n-m\right)}{B\left(\frac{g}{c}, \frac{r}{c}\right)}
	 &=\frac{\Gamma\left( (g + r) / c \right) }{\Gamma(g / c) \Gamma(r / c)}
	 \frac{\Gamma(g / c + m)}{\Gamma(m + 1)} 
	 \frac{\Gamma(r / c + n - m)}{\Gamma(n - m + 1)} 
	 \frac{\Gamma(n + 1)}{\Gamma\left((g + r) / c + n\right)} \\
\end{align*}
where, as $n, m \to \infty $,
\begin{align*}
	 \frac{\Gamma(g / c + m)}{\Gamma(m + 1)} 
	 \frac{\Gamma(r / c + n - m)}{\Gamma(n - m + 1)} 
	 \frac{\Gamma(n + 1)}{\Gamma\left((g + r) / c + n\right)} 
	&\approx  m^{g / c - 1} (n - m)^{r / c - 1} n^{- (g + r) / c + 1}  \\
	&\approx (x n)^{g / c - 1} \left( (1 - x) n \right) ^{r / c - 1} n^{- (g + r) / c + 1} \\
	&= x^{g / c - 1} (1 -x)^{r / c - 1} 
	\intertext{here the 
first $\approx$ is justified by the Stirling approximation $\frac{\Gamma(N + a)}{\Gamma(N + b)} \sim N^{a - b}$ when $N \to  \infty $, and the second $\approx$ is justified by the convergence of $X_n$ to $X_\infty  = x$.
}
.\end{align*}
Therefore,
\begin{lemma}
The probability density of $X_\infty $ is 
\[
	p(x) :=
	= \frac{\Gamma\left( (g + r) / c \right) }{\Gamma(g / c) \Gamma(r / c)}x^{g / c- 1} (1 -x)^{r / c - 1}
.\] 
\end{lemma}

\subsection{Exchangeability}
Polya's urn process is closely related to the concept of exchangeability. Exchangeable sequences of random variables have the property that the joint distribution of any finite subsequence remains unchanged when the order of elements is shuffled. Polya's urn is exchangeable because the order in which balls are drawn and replaced does not affect the distribution of colors.

 $\mathcal{E} = \cap \mathcal{E}_n$, where $\mathcal{E}_n$ is the $\sigma$-algebra generated by events invariant under all permutations of the first $n$ draws.

\begin{theorem}[de Finetti's Theorem]
	 If $X_1, X_2, \ldots$ are exchangeable then conditional on $\mathcal{E}, X_1, X_2, \ldots$ are independent and identically distributed.
\end{theorem}

The Polya's Urn is an exchangeable process, so the theorem tells us that the draws from the urn  $Y_i$ can be sampled independently from the same distribution (Bernoulli distribution) with Beta parameter. To write down the formula, we first let $D_i = 1$ if the $i$-th draw is green, and $D_i = 0$ if the $i$-th draw is red. Now
\begin{proposition}
	
\[
	P\left( D_1 = d_1, \ldots, D_n = d_n \right) 
	= \int_0^\infty  \prod_{i \le n} x^d_i (1 - x)^{1 - d_i} p(x) dx
.\] 
\end{proposition}



\subsection{Edge Reinforced Random Walk}

The Edge Reinforced Random Walk (ERRW) $(Y_k)_{k \ge 0}$ is a nearest-neighbor random walk in $\mathbb{Z}$ such that the transition probabilities depend completely on the number of times the two edges beside the walker have been visited. More precisely, let $L(\{z, z+1\} , k)$ denote the number of times the edge connecting sites $z$ and $z+1$ has been crossed by the time $k$.

Now, to define transition probabilities, first define the \textit{weight function} of each edge:
\[
	w(z, k) = 1 + L(\{z, z+1\} , k)
.\] 

Therefore the weight is linearly reinforced as the edge is being crossed. Now the transition probabilities are given by the weights:

\begin{align*}
	P(Y_{k+1} = z+1 | Y_k = z) &= \frac{w(z, k)}{w(z-1, k) + w(z, k)}\\
	P(Y_{k+1} = z-1 | Y_k = z) &= \frac{w(z-1, k)}{w(z-1, k) + w(z, k)}\\
.\end{align*}

Now, let $\lambda_{z, i}$, $i \ge 0$ denote the stopping time when the walk visits $z$ for the $i+1$-th time. We remark that $\lambda_{z, i}$ is always finite since the walk is known to be recurrent.

The local behavior of the walk, i.e. the evolution of weights of neighboring edges of $z$, is exactly a Polya's Urn:

\begin{proposition}
	Fix a site $z \in \mathbb{Z}$. Then:
	\begin{enumerate}
		\item 
			The process $\left(w(z-1, \lambda_{z, i}), w(z, \lambda_{z, i})\right)_{i \ge 0}$ is Markov.
		\item The processes associated with different sites are independent.

		\item 
	Then there is a Polya's Urn process $(R_i, G_i)_{i \ge 0}$ such that
	\[
		\left(w(z-1, \lambda_{z, i}), w(z, \lambda_{z, i})\right)_{i \ge 0}
		=_d
		\left( R_i, G_i \right) _{i \ge 0}
	.\]
	with parameter $c = 2$ and initial conditions are given by
	\end{enumerate}
	\begin{align*}
		(r, g) &= \begin{cases}
			(1,2) & z < 0\\
			(1,1) & z = 0\\
			(2,1) & z > 0
		\end{cases}
	.\end{align*}
\end{proposition}

\begin{remark}
	
Using this result, the \textit{drift} accumulated by the walk at the site $z$ 
\[
	\Delta_z^{(z, m)} := \sum_{k = 0}^{\lambda_{z, m}} E\left( Y_{k + 1} - Y_k | Y_1, \ldots, Y_k \right) \mathbf{1}\left( Y_k = z \right) 
\] 
can be computed explicitly from the Polya's Urn process corresponding to site $z$. When we have weight functions that are more general but at least still strictly increasing or decreasing, the same approach can work by generalizing the transition rules of P\'olya's urn a bit.

Asymptotic estimates of drifts are often useful in proving results about random walks. For example, this paper (Kosygina et al. 2022) establishes the convergence of a self-interacting random walk with weight function $w(n)^{-1} \sim  1 + c n^{-p} + o(n^{- 1 - \kappa})$ to Brownian Motion Perturbed at Extrema. For $p > \frac{1}{2}$, it is shown that $\Delta_z$ converges to $\text{sgn}(z) \gamma$ for some deterministic parameter $\gamma$. This allows us to understand behavior of the walk by ``summing up the drifts''.



\end{remark}

\subsection{References}

Durrett, Richard. Probability: Theory and Examples. 4th ed, Cambridge University Press, 2010.

Kosygina, Elena, et al. Convergence and Non-Convergence of Scaled Self-Interacting Random Walks to Brownian Motion Perturbed at Extrema. arXiv:2208.02589, arXiv, 4 Aug. 2022. arXiv.org, https://doi.org/10.48550/arXiv.2208.02589.

